\documentclass[PICOReport.tex]{subfiles}

%\newcolumntype{L}[1]{>{\raggedright\let\newline\\\arraybackslash\hspace{0pt}}m{#1}}
%\newcolumntype{K}[1]{>{\raggedright\centering\arraybackslash}m{#1}}

\begin{document}

In developing the theoretical framework for measuring the faint CMB
polarization signal, the CMB community has long recognized the impacts
of systematic errors due to instrumentation and data analysis.   A rich literature investigates the types of
systematic errors that confound the polarization measurement and ways
to mitigate them\cite{hu03}. 
Each on-orbit measurement of CMB polarization was limited by
sysetmatic errors until an in-depth study of the systematics was
performed and the post-processing data analysis removed them as well
as possible\cite{Bennett13,planck2016_xlvi}. 
Additionally, recently proposed CMB missions, such as LiteBird and \core, have placed
systematic error mitigation at the forefront of the case for their
mission\cite{hazumi2012,wallis2017,Natoli2018}.

\textbf{Note: Calculate a defensible level of acceptable systematics.
For example, integrate r $10^{-3}$ and $10^{-4}$ power from ell 2-10, 10-50,(do a signal to
noise calculation), then state what map rms that corresponds to.  }

A CMB mission aiming for the unprescedented sensitivity of PICO must control
systematic errors to avoid bias or an increased variance of the
science measurement.   
Systematics must be controlled or corrected to a level that enables the PICO
science goals (better than 1 nanoKelvin (\textbf{Note: make sure this is a
  correct statement - also be careful about reducing a complicated
  concept such as systematic errors to a single map rms number, could also
  quote this as a $\delta$r error}) in the map).
Mitigation of systematic errors is the most important reason (along
with the availability of broad wavelength coverage) to perform a
measurement of the CMB polarization from a space telescope; Compared
with a ground-based, sub-orbital, or even a space mission in 
low-Earth orbit, the L2 environment offers excellent stability as well
as the ability to observe large fractions of the sky on many time
scales without interference from the Sun, Earth, or Moon.
This redundancy of observations allows the checking of consisitency of
results and an improved ability to correct systematic errors in
post-processing analysis.

During the course of the PICO Study, a systematics working group
examined systematic errors affecting PICO,
Most  systematic errors can be mitigated by careful design and engineering of 
the spacecraft and instrument, and the use of present-day state-of-the-art
technology and data analysis tools.   
However, some systematic errors may
limit the precision of the B-mode measurement and the group studied
these in further detail. 
The work was based on the experience of the group's involvement with
past missions, in particular Planck, and in recent detailed studies of 
the \core\ and LiteBird  concepts.

End-to-end simulation of the experiment is an essential tool,
including realistic instabilities and non-idealities of the spacecraft,
telescope, instrument and folding in data post-processing techniques
used to mitigate the effects.  Systematics are coupled with the
spacecraft scan strategy, and the details of the 
data analysis pipeline.  During the study, the PICO team used 
 simulation and analysis tools developed for the Planck mission\cite{plank2015_xii_focalplane} and 
the \core\ mission concept, adapting them for PICO.  These tools allowed
a deeper examination of several key systematic errors. 

\subsubsection{List of Systematics}
The systematic errors face by PICO can be categorized into three broad categories 
1) Intensity-to-polarization leakage, 2) stability, and 3)
straylight.    These were prioritized for further study based on the
team's assessment of how well these systematics are understood by the
community, whether mitigation techniques exist - either in instrument
design or in data analysis.

\begin{table}[h!]
\centering
\scriptsize
 \begin{tabular}{p{3.8cm} p{4.2cm} p{4.2cm} p{4.2cm}}
 \hline
\textbf{Name} & \textbf{Description} & \textbf{State-of-the-art} & \textbf{Additional Possible Mitigation} \\
 \hline
\textbf{Leakage}& &\\
 Bandpass Mismatch\dotfill& 
 Edges and shapes of the the spectral filters vary from detector to detector. leaks T $\to$ P, P $\to$ P if the source's bandpass differs from calibrator's bandpass\cite{Hoang_2017} & Precise bandpass measurement\cite{Pajot_2010};
SRoll algorithm\cite{Planck_Lowell}; filtering technique\cite{CORE_systematics};   &
polarization modulation; full I/Q/U maps for individual detectors
                                                                                     mitigates;
                                                                                     additional
                                                                                     component
                                                                                     solution
                                                                                     (see
                                                                                     Banerji\&
                                                                                              Delabrouille
                                                                                               (in
                                                                                             prep).Current
                                                                                     techniques may be adequate  \\
Beam mismatch\dotfill& 
Beam shapes differ between detectors that are combined to reconstruct polarization; leaks T $\to$ P, P $\to$ P
& See Sect.~\ref{sec:angle} & Current techniques may be adequate.\\
Gain mismatch\dotfill&
Relative gain between detectors that are combined to reconstruct polarization; error leaks T$\to$P &
mission-average relative calibration demonstrated to 10$^{-4}$ to 10$^{-5}$ level \cite{Planck_Lowell}
&
Sect.~\ref{sec:gain} describes effects of stability in time in relative gains.  
\\
Time Response Accuracy and Stability\dotfill&
Uncertainty of detector in time constants (measurement errors, time variability) biases polarization angle, pointing and beam size. In a constant spin-rate mission (\pico) is degenerate with the beam shape. leaks T$\to$P, P$\to$P&
On-orbit reconstruction of time response to 0.1\% across a wide signal band\cite{planck2013_vii}, residuals corrected as part of beam and map-making algorithm\cite{Planck_Lowell}.
& Treatment of residual time response as part of the beam may be adequate.
\\
Readout Cross-talk\dotfill& 
Power in one detector leaks into other detectors
&
\planck's high-impedence bolometers with crosstalk measured at the level of 10$^{-3}$ did not impact CMB polarization science\cite{Planck_Lowell}.  Cross-talk low-impedence bolometers measured at XXX.
&
State-of-the-art (reconstruction and correction at map-making level) is adequate.
\\
Polarization Angle\dotfill& 
Uncertainty in polarization calibration leaks E$\to$B.
& 
Knowledge of astrophysical calibrators to 0.3$^\circ$\citet{Aumont+2018}; ground measurement to 0.9$^\circ$ reconstruction to 0.2$^\circ$ using $TB$ and $EB$ demonstrated by \planck\citet{Planck_Lowell}
& 
Polarization modulation; See Sect.~\ref{sec:angle} for discussion; Must mitigate through developing further analysis techniques.
\\

Cross-polarization\dotfill&
Q$\to$U rotation by the optical elements of the instrument.
&
Degenerate with polarization gain calibration.
&
Polarization modulation; State-of-the-art polarization gain calibration may be adequate.
\\
Chromatic beam shape\dotfill&
Beam shape is a function of source SED: measured using a planet, used to build a window function to correct CMB power spectrum.
&
\planck\ simulations and parameterization as part of the likelihood.
&
Should be further investigated in Phase A of a mission using physical optics simulations.  
\\
\hline 
\textbf{Stability} & & \\

Pointing jitter\dotfill&
Random pointing error mixes T, E and B at small angular scale
&
Pointing reconstruction in \planck\ to 
&
\\

Gain Stability\dotfill& 
Time-variation of detector gain due to time variability of bath temperature variations, optical power.
& 
Reconstruction of time variability of gain to 0.2\% in \planck\cite{Planck_Lowell}.
&
See Sect.~\ref{sec:gain}; Gain fluctuations in \pico\ on the level of XXX\% on time scales of YYY can be corrected in post-processing.
\\
\hline
\textbf{Straylight}& & \\
Far Sidelobes\dotfill& 
Pickup of Galactic signals at large angles from the main beam axis; Spillover can be highly polarized.
& 
\planck\ validated straylight model in anechoic chamber to -80~dBi\cite{Tauber2010}.
&
Design of optical system and baffling, informed by telescope straylight simulations. See Sect.~\ref{sec:fsl} for a study of beams calculated with a physical optics code for the \pico\ telescope and simulated Galactic pickup during the reference mission.\\
 \hline
\textbf{Other} \\
Residual correlated cosmic ray hits\dotfill&
detectors experience correlated cosmic ray hits below detection threshold resulting in misestimated noise covariance.
&
\planck/HFI found the 5\% percent noise correlation due to this effect did not impact results\cite{Planck_Lowell}. 
&
Detector design to reduce cosmic ray cross-section; Current analysis techniques (accounting for correlated noise) may be adequate.
\\
\hline
 \end{tabular}
\caption{\label{tbl:SystematicsList} Systematic errors expected in \pico's measurement of CMB polarization, with assessment of currently known mitigation techniques.  Those systematic errors found to be most likely to impact \pico\ are described further in the text.}
 \end{table}

\subsubsection{Absolute polarization angle calibration}
\label{sec:angle}

The rotation of the CMB polarization can have different causes,
including 1. a birefringent primordial Universe, or a Faraday rotation
due a primordial magnetic field \citep{Pogosian+2018}, 2. birefringent
foregrounds, or interaction with the Galactic magnetic field,
3. systematic effects in the instrument, and in particular an error on
the actual direction of polarization measured by each detector.  
While the first two sources create a rotation that may depend on scale,
position and/or frequency, the latter is expected to mostly depend on
the detector considered. 

A rotation {\prang} of the direction of polarization mixes the $Q$ and $U$ Stokes parameters via
$Q\pm iU \longrightarrow e^{\mp i 2 \prang} (Q\pm iU)$
and affects the power spectra via (assuming the rotation to be independent on scale and location)
\begin{subequations}
\begin{align}
C^{TT}_\ell &\longrightarrow & C^{TT}_\ell                                             &= & C^{TT}_\ell \\
C^{TE}_\ell &\longrightarrow & \cos 2\prang\  C^{TE}_\ell                                &\sim & \left(1 - 2\prang^2\right)\ C^{TE}_\ell \\
C^{EE}_\ell &\longrightarrow & \cos^2 2\prang\  C^{EE}_\ell + \sin^2 2\prang\  C^{BB}_\ell &\sim & C^{EE}_\ell - 4\prang^2\ \left(C^{EE}_\ell - C^{BB}_\ell\right) \\
C^{BB}_\ell &\longrightarrow & \sin^2 2\prang\  C^{EE}_\ell + \cos^2 2\prang\  C^{BB}_\ell &\sim & C^{BB}_\ell + 4\prang^2\ \left(C^{EE}_\ell - C^{BB}_\ell\right)\\
C^{TB}_\ell &\longrightarrow & \sin 2\prang\  C^{TE}_\ell                                &\sim & 2\prang\  C^{TE}_\ell \\
C^{EB}_\ell &\longrightarrow & \sin 2\prang \cos 2\prang \left(C^{EE}_\ell -  C^{BB}_\ell\right)  &\sim & 2\prang\ \left(C^{EE}_\ell -  C^{BB}_\ell\right)
\end{align}
\end{subequations}
as illustrated in Fig.~\ref{fig:rot_bb_tb_eb}.

%------------------------------------------------------------------------------------------
\begin{figure}[htb]
\includegraphics[width=0.8\textwidth]{images/PICO_rotate_eb2_v0.\suffix}
\caption{\label{fig:rot_bb_tb_eb} Effect of a rotation of the angle of polarization, assuming the Planck 2018 $\Lambda$-CDM best fit model \citep{Planck2018_VI} and expected \pico performances, with a perfect delensing \EFH{(the beam related systematics effects are still arbitrary; to be improved or removed)}.}
\end{figure}
%------------------------------------------------------------------------------------------
%
%------------------------------------------------------------------------------------------
\begin{figure}[htb]
\includegraphics[width=0.40\textwidth]{images/PICO_sens2_v0_F5p0_f0p5_n0p62_k4_a1p0_2_4000.\suffix}
\includegraphics[width=0.40\textwidth]{images/PICO_sens2_v0_F15p0_f0p5_n0p62_k4_a1p0_2_4000.\suffix}
\caption{\label{fig:rot_sens_0} Upper panels: signal to noise ratio of the polarization angle {\prang} measurement
by $EB$ (blue lines), $TB$ (green lines) and $BB$ (red lines), assuming either no delensing (solid lines) 
or perfect delensing (dashes); the shaded area is $|\prang|/\sigma_\prang < 3$.
Lower panels: degradation on measurement of $r$, for $r=10^{-2},\ 10^{-3},\ 10^{-4}$ (magenta, orange and cyan lines, respectively),
either with no delensing (solid lines) or perfect delensing (dashes).
The underlying cosmology is Planck 2018 $\Lambda$-CDM model (with $\tau = 0.054$), and assuming a polarized noise of rms = $0.62 \mu K.\arcmin$ and power spectrum $(1 + (\ell_{\rm knee}/\ell)^n)$ with $\ell_{\rm knee}=4$ and $n=1$, with the analysis done on the multipole range $[2,4000]$ over a sky fraction $\fsky=0.5$. The beam FWHM$=5\arcmin$ on the \emph{lhs} and $15\arcmin$
on the \emph{rhs} panels.}
\end{figure}
%
% \begin{figure*}[htb]
% \includegraphics[width=0.5\textwidth]{fig_efh/PICO_sens2_v0_F15.0_f0.5_n0.62_k4_a1.0_2_4000.\suffix}
% \caption{\label{fig:rot_sens_1} Same as Fig.~\ref{fig:rot_sens_0}, with a FWHM=$15\arcmin$.}
% \end{figure*}
%
\begin{figure}[htb]
\includegraphics[width=0.40\textwidth]{images/PICO_sens2_v0_F5p0_f0p5_n0p62_k4_a1p0_20_4000.\suffix}
\includegraphics[width=0.40\textwidth]{images/PICO_sens2_v0_F5p0_f0p5_n1p86_k4_a1p0_2_4000.\suffix}
\caption{\label{fig:rot_sens_2} Same as Fig.~\ref{fig:rot_sens_0}, left panels, reducing the multipole range $[20,4000]$ (\emph{lhs}) or with a noise rms multiplied by 3 (\emph{rhs}).}
\end{figure}
% %
% \begin{figure*}[htb]
% \caption{\label{fig:rot_sens_3} Same as Fig.~\ref{fig:rot_sens_0}, with a noise rms multiplied by 3.}
% \end{figure*}
% %------------------------------------------------------------------------------------------

In Planck, the ground measurements of the detectors orientation had an error of $\pm 0.9\degree$ (rel.) $\pm 0.3\degree$ (abs.) \citep{Rosset+2010}.

The most recent constraints on cosmological birefringence (or systematic rotation) was set in \citet{Planck2016_XLIX}, looking for residual signal in $TB$ and $EB$ spectra, but are dominated by the uncertainties on the detector orientations.

In \pico, the relative rotation of the detectors, could be measured with a good accuracy (a few $0.1\arcmin$ ?, \EFH{refs?}) on the CMB, but the overall rotation is difficult to determine.
Known polarized sources, such as the Crab Nebula, could be used to do that but \citet{Aumont+2018} show that the current uncertainty of $0.33\degree = 20\arcmin$ on the Crab polarization orientation, obtained when combining all the available measurements, 
would not the measurement of tensorial $B$ modes below $r \sim 0.01$ (assuming everything else to be nominal), far from \pico's target.

Figures \ref{fig:rot_sens_0} and \ref{fig:rot_sens_2} show how the measurement of $r$ by \pico\ is degraded because of an overall rotation of polarization, and how $TB$ and $EB$ can be used to monitor precisely this rotation, assuming that the only source of polarization rotation is instrumental.
These results are obtained assuming the spectra to have a Gaussian likelihood, with a variance $\propto 1/\fsky$, and ignoring the foreground contributions.



$TB$ and $EB$ spectra can detect and measure a global polarisation rotation at levels ($~0.1 \arcmin$) well below those affecting $r$ measurements in $BB$ ($> 1 \arcmin$).
However, in doing such studies, one can not ignore other important aspects of the measurements of CMB polarization:
For instance, one has to consider how the delensing will affect the $B$ maps and their final noise levels.
At large scales, the interaction with foregrounds makes masking necessary, which creates some extra $E$ to $B$ leakage, and complicates the estimation of $B$ modes and of their likelihood. It also necessitates a proper component separation technique, especially in presence of finite and probably mismatched spectral bandwith of the detectors.
The time correlation of the noise (with a $1/f$ spectral shape), will also complicate the analysis.
At intermediate and small scales, many other systematics, like those related to beams, will also have to be modeled.



\subsubsection{Gain Stability}
\label{sec:gain}
Photometric calibration is the process of converting the raw output of the receivers into a physically-meaningful quantity, such as thermodynamic temperature or brightness. As CMB receivers are usually linear, this process reduces to the characterization of the \emph{gain factor} $G$:
\begin{equation}
\label{eq:gainStability:Response}
y(t) = G(t) \times T\bigl(\vec{x}(t)\bigr) + n(t),
\end{equation}
where $y(t)$ is the timestream of raw samples produced by the detector, $G(t)$ is the gain factor (which we allow to vary with time), $T$ is the sky temperature observed along direction $\vec{x}(t)$ (which varies with time as the spacecraft spins), and $n(t)$ is a noise term that includes both uncorrelated and correlated noise. It is assumed that the timescale of variation in $G$ ($G / \dot{G}$) is much longer than the typical timescale of variations in $T$: in the case of Planck, this was of the order of several days.

In the case of space CMB experiments, the characterization of $G(t)$ is commonly done using the signal caused by the motion of the spacecraft with respect to the rest frame of the CMB itself. This signal is commonly called the \emph{dipole}, as its most significant contribute is at multipole $\ell=1$. For the PICO concept study, we evaluated the impact of noise in the estimation of $G(t)$ using the tools developed for the Planck/LFI instrument and the CORE mission proposal. The quality of the estimate depends on the noise level of the receiverse, but also on the details of the scanning strategy. The Planck/LFI experiment, because of a poor choice of the scanning strategy parameters (namely, a too slow precession motion), was forced to avoid using one year out of four in the 2015 data release [REF]. We can anticipate that this problem is not expected in PICO, thanks to the significantly faster precession envisaged.

In order to test the impact of calibration uncertainties, we have run the following analysis:
\begin{enumerate}
\item\label{step:gainStability:perfectObservation} We simulated the observation of the sky, assuming four receivers and the nominal scanning strategy. We included both white noise and $1/f$ noise. The sky only contained CMB anisotropies, plus the CMB dipole.
\item\label{step:gainStability:calibration} We ran the calibration code to fit the dipole against the raw data simulated during step~\ref{step:gainStability:perfectObservation}.
\item\label{step:gainStability:imperfectObservation} We simulated again the observation of the sky, but this time we used the values of $G$ computed during step~\ref{step:gainStability:calibration}, which contain errors due to the presence of $n(t)$ and the CMB signal in Eq.~\eqref{eq:gainStability:Response}. The noise in the output map is therefore the sum of the noise in the error on $G$ and the term $n$.
\end{enumerate}
The presence of foregrounds in the sky signal would cause a bias in the estimation of the calibration constants, due to the presence of large scale features in the Milky Way at microwave frequencies. A full data analysis pipeline for PICO should pair the calibration step with the component separation step, following a schema similar to what has been done by the Planck/LFI team for the 2018 data release [CITATION]: the application of the calibration code should be followed by a component separation analysis, and these two steps should be iterated until the result converge to a solution. In this analysis we assume to study the calibration at the last iteration, when the components have already been properly separated.

Results of the simulation are shown in Figures~XXX and YYY. The scanning strategy employed by PICO allows for a much better calibration than in the case of Planck's, thanks to the much faster precession.

\subsubsection{Far Sidelobe Pickup}
\label{sec:fsl}
The main beam (within a few degrees of the axis of beam response) in a CMB mission can be measured to high precision using the planets as compact sources.  Measurement of each detector's response to signals more than a few degrees off axis, which tends to be at a very low level (more than --80dB less than the peak response) but spread over a very large solid angle, is difficult to do pre-launch, and may not be done accurately after launch.  Nonetheless, this far sidelobe can couple bright Galactic signal from many tens of degrees off-axis and confuse it with polarized signal from the CMB off the Galactic plane.  

To evaluate this systematic error, GRASP software\footnote{https://www.ticra.com} was used to compute the \pico\ telescope's pickup over the full sky.  This full-sky beam was convolved with a polarized Galactic signal and a full \pico\ mission scan using the simulation pipeline.  The far sidelibe pickup was estimated to contribute less than XXX to the B-mode angular power spectrum and thus an error in $r$ of YYY.

In a real mission due to the difficulties of measuring this beam, physical optics simulation capabilities must be maintained and validated as well as possible with on-orbit data.

\subsubsection{Key Findings}
Understanding and controlling the effects of systematic errors in a
next-generation CMB probe is critical.

The raw sensitivity of the instrument should include enough margin
that data subsets can independently archieve the science goals.
This allows testing of the results in the data analysis and additional
data cuts, if needed.

In a PICO mission's phase A, a complete end-to-end system-level
simulation software facility would be developed to assist the team in setting 
requirements and conducting trades between subsystem requirements while
realistically accounting for post-processing mitigation.  Any future
CMB mission is likely to have similar orbit  
and scan characteristics to those of PICO, thus there is an opportunity for NASA and
the CMB community to invest in further development of this capability now.


\end{document}

%\begin{figure}[!htb]
%\centering
%\includegraphics[width=4cm]{images/example}
%\caption{example}
%\label{fig:im_3}
%\end{figure}
